\documentclass{beamer}

% \usepackage[margin = 0.7in]{geometry}
\usepackage{tikz}
\usepackage{titlesec}
\usepackage{siunitx}
\usepackage[table]{xcolor}
\usepackage{ifthen}
\usepackage{nicematrix}
\usepackage{booktabs}
\usepackage{tabularx}

\usetikzlibrary{calc}
\usetikzlibrary{positioning}

\renewcommand{\tabularxcolumn}[1]{m{#1}}

\definecolor{colorRed} {HTML} {E44E54}
\colorlet{colorLightRed}{colorRed!70!white}

\colorlet{colorPill}{black!30!white}

\definecolor{colorPastelGreen} {HTML} {97B3AE}

\colorlet{colorNonSun}{white!15!colorPastelGreen}
\colorlet{colorNonSunDark}{black!30!colorNonSun}

\colorlet{colorSun}{colorLightRed}
\colorlet{colorSunDark}{black!30!colorSun}

\definecolor{colorP1} {HTML} {D2E0D3}
\definecolor{colorP2} {HTML} {F0DDD6}
\definecolor{colorP3} {HTML} {F2C3B9}
\definecolor{colorP4} {HTML} {D6CBBF}

\colorlet{colorP1Dark}{black!60!colorP1}
\colorlet{colorP2Dark}{black!60!colorP2}
\colorlet{colorP3Dark}{black!60!colorP3}
\colorlet{colorP4Dark}{black!60!colorP4}

\titleformat{\chapter} [display]
{\bfseries\normalfont\huge\filright\sffamily\vspace{-2cm}}
{\color{colorNonSunDark} \Large\texttt{WEEK {01} \hfill OCT 05 -- OCT 11} \vspace{1em}}
{1pc}
{\color{colorNonSunDark}\titlerule\vspace{1.0em}\scshape\centering\color{colorNonSunDark}}
[\color{colorNonSunDark}\vspace{0.65em}{\titlerule[1pt]}]

\titleformat{\section}
{\vspace{-0.35em}\bfseries\normalfont\large\filright\sffamily}
{\thesection}
{8pt}
{\scshape}
{}

\titleformat{\subsection}
{\vspace{-0.35em}\bfseries\normalfont\filright\sffamily}
% {}
% {0pt}
{\thesubsection}
{8pt}
{\scshape}
{}


% \usepackage{minted}

\usepackage[size=custom,width=16,height=9,scale=0.35]{beamerposter}

\usepackage[T1]{fontenc}
% \usepackage[a4paper, margin = 1in]{geometry}
\usepackage{tcolorbox}
\usepackage{graphicx}
\usepackage{lmodern}
\usepackage{tabularx}
% \usepackage{ebgaramond}
\usepackage{tikz}
\usepackage{blindtext}
\usepackage{tcolorbox}
\usepackage{svg}
% \usepackage{fontspec}
\usepackage{fontspec}
\usepackage{booktabs}

\usepackage[style = ieee]{biblatex}

% \setmonofont{Vaulto-ExtraBold.ttf}

\usetikzlibrary{calc}

% \usefonttheme{serif}
% \usepackage{beamercolorthemedracula}

\addbibresource{abstract.bib} %Import the bibliography file

\usetheme{metropolis}
\usecolortheme{seedinc}

\definecolor{forestbg}      {RGB} {31, 30, 24}

\title{Seed Incubation Plant}
\date{\today}
\def\author{
    19. Anjana Roy \hfill KTE22EC020 \\
    \vspace{3pt}
    25. Aswatheertha T T \hfill KTE22EC026 \\
    \vspace{3pt}
    26. Athira Madhusoodanan \hfill KTE22EC027 \\
    \vspace{3pt}
    28. Daniel V Mathew \hfill KTE22EC029 \\
}

\def\supervisor{
    Prof. David Solomon George
}

\def\leafColorPrimaryName{Sage}
\def\leafColorPrimaryRGB{138, 184, 114}
\def\leafColorPrimaryHex{\#8AB872}
\definecolor{leafColorPrimary} {RGB} {\leafColorPrimaryRGB}
%%
\def\leafColorSecondaryName{Cactus Spike}
\def\leafColorSecondaryRGB{204, 224, 153}
\def\leafColorSecondaryHex{\#CCE099}
\definecolor{leafColorSecondary} {RGB} {\leafColorSecondaryRGB}


\newcommand\tableBox[2] {
    \begin{tcolorbox}[
            coltitle = leafColorPrimary,
            colbacktitle = leafColorSecondary,
            colback = leafColorSecondary,
            colframe = leafColorSecondary,
            title=#1,
            fonttitle=\bfseries,
            detach title,
        ]
        \begin{minipage}[t]{0.1\textwidth}
            \begin{flushleft}
                \tcbtitle
            \end{flushleft}
        \end{minipage}
        \begin{minipage}[t]{0.9\textwidth}
            #2
        \end{minipage}
    \end{tcolorbox}
}

\setbeamertemplate{background} {
    \begin{tikzpicture} [remember picture, overlay]
        \draw [opacity = 0.15, black]
        (current page.south east) ++(-3.5, 1)
        node (logo) [black, inner sep = 0pt, anchor = center] {
            \begin{minipage}{11cm}
                \includegraphics[width = 11cm]{icon_line.pdf}
            \end{minipage}
        }
        ;
    \end{tikzpicture}
}

% \institute{RIT, Kottayam}

\renewcommand{\maketitle} {

    \vfill

    \begin{tikzpicture} [remember picture, overlay]

        \draw
        (current page.north) ++(0, -1cm)
        node (title) [inner sep = 0pt, anchor = north] {

            \begin{minipage} {\textwidth}
                {
                    {\color{leafColorSecondary} \rule{\textwidth}{0.5pt}}
                    \vspace{3pt}

                    {\huge \scshape Seed Incubation Plant}

                    \vspace{3pt}
                    {\color{leafColorSecondary} \rule{\textwidth}{0.75pt}}
                }

            \end{minipage}

        }
        ;

        \draw

        (title.south)
        node (group) [inner sep = 0pt, anchor = north] {

            \begin{minipage} {\textwidth}
                {\scshape group 10}
            \end{minipage}

        }
        ;

        \draw

        (current page.south) ++(0, 1cm)

        node (something) [inner sep = 0pt, anchor = south] {

            \begin{minipage}[b] {\textwidth}

                \small
                Team Members
                \hfill
                Guide Name \& signature

                \vspace{8pt}

                \noindent
                {\begin{minipage}[b] {0.375\textwidth}
                    \vfill
                    \raggedright
                    {\linespread{1.2} \bfseries \author}
                \end{minipage}
                \hfill
                \begin{minipage}[b] {0.4\textwidth}
                    \vfill
                    \raggedleft
                    {\linespread{1.2} \bfseries \supervisor}
                \end{minipage}}

            \end{minipage}
        }

        ;

    \end{tikzpicture}

    \vfill

}

\def\objImageRadius{1.75cm}
\def\objImageWidth{2.625cm}

\newcommand\includeObjImage[4]{
    \begin{tikzpicture} [
            transform shape,
            scale = 1,
        ]

        \draw
        (0, 0) coordinate (botLeftCorner)
        (botLeftCorner) ++(#2, #2) coordinate (topRightCorner)

        (botLeftCorner |- topRightCorner) coordinate (topLeftCorner)
        (botLeftCorner -| topRightCorner) coordinate (botRightCorner)
        ;

        \filldraw [
            leafColorSecondary,
            rounded corners = #4,
        ]
        (botLeftCorner) rectangle (topRightCorner)
        % circle (\objImageRadius)
        ;

        \draw [
            ultra thick,
            leafColorPrimary,
            rounded corners = #4,
        ]
        (botLeftCorner) rectangle (topRightCorner)
        % circle (\objImageRadius)
        ;

        \draw

        ($(botLeftCorner)!0.5!(topRightCorner)$)
        node [
            anchor = center,
        ] {
            \includegraphics[width = #3]{#1}
        }
        ;

    \end{tikzpicture}

    \thispagestyle{empty}
}

\begin{document}

\maketitle

\begin{frame} {Introduction}

    {
        \begin{minipage} {0.415\textwidth}
            \begin{center}
                \includeObjImage{ext_right.png}{5}{4.5cm}{0.5cm}
            \end{center}
        \end{minipage}
        \begin{minipage} {0.55\textwidth}
            \textbf{Seed Incubation Plant} aims to optimize the germination and
            growth process of seeds by providing a controlled and
            monitored environment.
            While the \textbf{Environment Monitoring and Control Unit} (EMCU) part of the
            incubator takes care of environmental factors such as temperature,
            humidity, lighting, $\mbox{O}_2$, etc, the
            \textbf{Growth Monitoring Unit} (GMU) closely monitors and collects growth related
            information from each of the germinating saplings.
            By the help of \textbf{SINC}, a companion app, the status of the incubator
            can be closely watched.

        \end{minipage}
    }

\end{frame}

\begin{frame} {Objectives}

    \begin{center}
        \begin{tabularx} {\textwidth} {
                *{3}{>{\centering \arraybackslash}X}
            }

            \includeObjImage{enclosure.png}{4}{2.825cm}{0.5cm}
            &
            \includeObjImage{growth.png}{4}{2.625cm}{0.5cm}
            &
            \includeObjImage{sinc.png}{4}{3.025cm}{0.5cm}
            \\

            \begin{itemize}
                \item Implement the sensor and actuator sides of the EMCU, thus
                    driving it into a cohesive system.
            \end{itemize}
            &
            \begin{itemize}
                \item Implement an ESP32 CAM based Growth Monitoring Unit (GMU) powered by
                    a TinyML model.
            \end{itemize}
            &
            \begin{itemize}
                \item Take the development of SINC, a companion app, to the alpha
                    stage.
            \end{itemize}
            \\

        \end{tabularx}
    \end{center}


\end{frame}

\begin{frame} {Literature Review}

    \begin{tabularx} {\textwidth} {
            >{\bfseries}m{2.0cm}X
        }

        Title: & \textbf{Internet of Things (IoT) Based Greenhouse Monitoring and Controlling System Using ESP-32}
        \\

        Authors: & J. Seetaram, A. Bhavya, C. Tarun and V. Sameera
        \\

        DOI: & \texttt{10.17148/IJARCCE.2024.13605}

        \\

        Publisher: & \textit{International Journal of Advanced Research in Computer and Communication Engineering} \vfill
        \\

        Result: &
        The system successfully performed real-time monitoring and automated environmental control.
        Demonstrated fully autonomous and remote operation, reducing human intervention.

        \\

    \end{tabularx}

    {\color{leafColorSecondary} \rule{\textwidth}{0.75pt}}

    \textbf{Relevant Features \& Insights:}

    \begin{itemize}
        \item IoT-enabled remote control with real-time sensor data visuals.

        \item User-friendly UI designed on Blynk for farmers.

        \item Focus on sustainability, resource optimization, and precision agriculture.
    \end{itemize}

\end{frame}

\begin{frame} {Work Plan}

    \renewcommand\arraystretch{1.75}
    \renewcommand\midrule{{\color{leafColorSecondary} \vspace{-8pt}\rule{\textwidth}{0.75pt}\vspace{8pt}}}


    \begin{tabularx} {\textwidth} {
            >{\bfseries \raggedright \arraybackslash}m{2cm} >{\raggedright \arraybackslash}X
        }

        Week 01 - 02
        &
        Refinement of Enclosure design, data collection for TinyML model, and setting up
        the specification for the communication protocol and UI of SINC.
        \\ \midrule

        Week 03 - 04
        &
        Reimplement the Electronics Bay, including the connections that binds Multiplexer and
        Relay boards. Implement the outer frame of the Enclosure. Starts developing GMU and
        SINC app.
        \\ \midrule

        Week 05 - 06
        &
        Implement Thermal / Exhaust system along with the Top and Side Hatches. Testing
        the User Interface of SINC app.
        \\ \midrule

        Week 07 - 08
        &
        Implement Air Moisture and Lighting systems along with the Rail Mechanism for GMU. Testing
        of TinyML model on ESP32 CAM module. Testing of SINC app to its full specification.
        \\ \midrule

        Week 09
        &
        Integration of subsystems and further testing. Report preparation and documentation.

        \\

    \end{tabularx}

\end{frame}

\begin{frame} {References}

    \nocite{interEsp, ieeeComp, ieeeGreenAI, tinymlEnv}

    \nopagebreak

    \begingroup \let\clearpage\relax
    \printbibliography[heading = none] %Prints bibliography
    \endgroup

\end{frame}

\begin{frame} {}

    \begin{tikzpicture} [remember picture, overlay]

        \draw [
            leafColorSecondary,
            thick,
            line cap = round,
        ]
        (current page.center) ++(1.75, 0) coordinate (lineEnd)
        (current page.center) ++(-1.75, 0) coordinate (lineStart)
        (current page.center) coordinate (line)

        (lineStart) -- (lineEnd)
        ;

        \draw
        (line.north)
        node (logo) [inner sep = 0pt, anchor = south] {
            \begin{minipage} {10cm}
                \centering
                \includegraphics[width = 1cm]{icon_line.pdf}
                \vspace{10pt}
            \end{minipage}
        }
        ;

        \draw
        (line.south)
        node (ty) [inner sep = 0pt, anchor = north] {
            \begin{minipage} {10cm}
                \centering
                \vspace{10pt}
                \color{forestbg}
                \scshape
                Thank You

            \end{minipage}
        }
        ;

    \end{tikzpicture}

\end{frame}

\end{document}
